\begin{frame}[c]
   \vspace{0.56cm}
   \begin{columns}[c,totalwidth=1.05\textwidth]
      \begin{column}{0.6\textwidth}
         \hspace*{-0.24cm}%
         \begin{minipage}{\linewidth}
            \begin{itemize}
               \item Los grafos que vamos a considerar son \alert{localmente finitos} y \alert{sin vértices aislados}.
            \end{itemize}
         \end{minipage}
      \end{column}
      \begin{column}{0.45\textwidth}
         \hspace*{-0.24cm}%
         \begin{minipage}{\linewidth}
            \centering\small
            \setlength{\leftmargini}{1.2em}
            \uncover<2->{%
               \begin{tikzpicture}[
  scale=0.5,
  vnode/.style={circle, fill=FancyBlue!40, draw=black, thick, text=black,
                minimum size=6mm, inner sep=0pt, font=\footnotesize}
]
  % Grafo G_4 de la Figura 1.1.6: v_1 es un vértice aislado
  \node[vnode] (n1) at (2,4) {$v_1$};
  \node[vnode] (n2) at (4,3) {$v_2$};
  \node[vnode] (n3) at (4,1) {$v_3$};
  \node[vnode] (n4) at (2,0) {$v_4$};
  \node[vnode] (n5) at (0,1) {$v_5$};
  \node[vnode] (n6) at (0,3) {$v_6$};

  \foreach \from/\to in {n3/n4, n2/n3, n2/n5, n2/n6, n3/n5, n3/n6, n5/n6}
    \draw[thick] (\from) -- (\to);
\end{tikzpicture}
\par
               \vspace{0.15cm}
               \begin{minipage}{0.7\linewidth}
                  \begin{itemize}
                     \item $v_1$ es aislado
                  \end{itemize}
               \end{minipage}\par
            }
            \vspace{0.3cm}
            \uncover<3->{%
               \begin{tikzpicture}[
  scale=0.4,
  vnode/.style={circle, fill=FancyBlue!40, draw=black, thick, text=black,
                minimum size=6mm, inner sep=0pt, font=\footnotesize}
]
  % Grafo del Ejemplo 1.1.20: no es localmente finito, deg(v_1) no es finito
  \node[vnode] (n1) at (0,2.5) {$v_1$};
  \node[vnode] (n2) at (0,0)   {$v_2$};
  \node[vnode] (n3) at (2,0)   {$v_3$};
  \node[vnode] (n4) at (4,0)   {$v_4$};
  \node[vnode] (n5) at (6,0)   {$v_5$};

  \foreach \to in {n2, n3, n4, n5}
    \draw[thick] (n1) -- (\to);
  \draw[thick] (n1) -- (8.6,0);

  \node at (9.6,0) {$\cdots$};
\end{tikzpicture}
\par
               \vspace{0.15cm}
               \begin{minipage}{0.9\linewidth}
                  \begin{itemize}
                     \item No es localmente finito
                  \end{itemize}
               \end{minipage}
            }
         \end{minipage}
      \end{column}
   \end{columns}
\end{frame}
\begin{frame}[c]
   \vspace*{0.83cm}% centrado respecto al área bajo la cabecera
   \begin{columns}[c,totalwidth=1.05\textwidth]
      \begin{column}{0.6\textwidth}
         \hfill
         \begin{minipage}{0.9\linewidth}
            \begin{Def}
               Sea $G=(V,E)$ un grafo. Para cada $v \in V$, denotamos por $A_v$ el conjunto de todos los vértices adyacentes a $v$.\vspace{0.3cm}
            \end{Def}
         \end{minipage}
         \hfill\null
      \end{column}
      \begin{column}{0.45\textwidth}
         \centering
         \uncover<2->{\begin{tikzpicture}[
  vnode/.style={circle, fill=FancyBlue!40, draw=black, thick, text=black,
                minimum size=6mm, inner sep=0pt, font=\footnotesize}
]
  % Vértices: v2 en el centro
  \coordinate (v1) at (90:1.25);
  \coordinate (v2) at (0,0);
  \coordinate (v3) at (-45:1.25);
  \coordinate (v4) at (-90:1.25);
  \coordinate (v5) at (225:1.25);

  % Aristas: v1v2, v2v3, v2v4, v2v5
  \foreach \i in {1,3,4,5}
    \draw[thick] (v2) -- (v\i);

  \foreach \i in {1,...,5}
    \node[vnode] at (v\i) {$v_{\i}$};
\end{tikzpicture}
}

         \vspace{0.4cm}
         \begin{minipage}{0.95\linewidth}
            \begin{itemize}
               \item<3-> $A_{v_1}=\{v_2\}$
               \item<4-> $A_{v_2}=\{v_1,v_3,v_4,v_5\}$
            \end{itemize}
         \end{minipage}
      \end{column}
   \end{columns}
\end{frame}

\begin{frame}[c]
   \vspace*{0.76cm}% centrado respecto al área bajo la cabecera
   \hspace*{0.42cm}% equilibra las brechas laterales
   \begin{minipage}{\textwidth}
   \begin{columns}[c]
      \begin{column}{0.48\textwidth}
         \begin{prop}
            El conjunto $\{A_v \mid v \in V\}$ forma una \alert{subbase} para una topología sobre $V$.\vspace{0.3cm}
         \end{prop}
         \vspace{\baselineskip}
         \begin{itemize}
            \item<2-> A los $A_v$ los llamamos \alert{abiertos generadores}.
         \end{itemize}
      \end{column}
      \begin{column}{0.48\textwidth}
         \centering
         \uncover<3->{Abiertos generadores}\\[0.1cm]
         \uncover<4->{\includegraphics[height=0.55cm,angle=45]{Images/Flecha_azul.pdf}\\[0.1cm]
            Intersecciones finitas}\\[0.1cm]
         \uncover<5->{\includegraphics[height=0.55cm,angle=45]{Images/Flecha_azul.pdf}\\[0.1cm]
            Uniones arbitrarias}\\[0.1cm]
         \uncover<6->{\includegraphics[height=0.55cm,angle=45]{Images/Flecha_azul.pdf}\\[0.1cm]
            Topología}
      \end{column}
   \end{columns}
   \end{minipage}
\end{frame}

\begin{frame}[c]
   \vspace*{0.5cm}% centrado respecto al área bajo la cabecera
   \begin{prop}
      La topología generada por $\{A_v \mid v \in V\}$ es de \alert{Alexandroff}. La llamamos \alert{topología de adyacencia} sobre $G$ y la denotamos por $\tau_{\mathcal{A},G}$.\vspace{0.3cm}
   \end{prop}
   \uncover<2->{%
   \begin{prop}
      Para cada $x \in V$ tenemos \alert{$U_x=\bigcap_{y\in A_x}A_y$}.\vspace{0.3cm}
   \end{prop}}
\end{frame}

\begin{frame}[c]
   \vspace*{1.07cm}% centrado respecto al área bajo la cabecera
   \hspace*{0.47cm}% equilibra las brechas laterales
   \begin{minipage}{\textwidth}
   \begin{columns}[c]
      \begin{column}{0.3\textwidth}
         \centering
         \begin{tikzpicture}[
  vnode/.style={circle, fill=FancyBlue!40, draw=black, thick, text=black,
                minimum size=6mm, inner sep=0pt, font=\footnotesize}
]
  % Vértices: v2 en el centro
  \coordinate (v1) at (90:1.25);
  \coordinate (v2) at (0,0);
  \coordinate (v3) at (-45:1.25);
  \coordinate (v4) at (-90:1.25);
  \coordinate (v5) at (225:1.25);

  % Aristas: v1v2, v2v3, v2v4, v2v5
  \foreach \i in {1,3,4,5}
    \draw[thick] (v2) -- (v\i);

  \foreach \i in {1,...,5}
    \node[vnode] at (v\i) {$v_{\i}$};
\end{tikzpicture}

      \end{column}
      \begin{column}{0.68\textwidth}
         \begin{itemize}
            \item[] Por ejemplo:
            \item $U_{v_1} = A_{v_2} = \{v_1,v_3,v_4,v_5\}$
            \item<2-> $U_{v_2} = A_{v_1}\cap A_{v_3}\cap A_{v_4}\cap A_{v_5}$
            \item<2->[] $\phantom{U_{v_2}} = \{v_2\}\cap\{v_2\}\cap\{v_2\}\cap\{v_2\}$
            \item<2->[] $\phantom{U_{v_2}} = \{v_2\}$
         \end{itemize}
      \end{column}
   \end{columns}
   \end{minipage}
\end{frame}

\begin{frame}[c]
   \vspace*{0.32cm}% centrado respecto al área bajo la cabecera
   \centering
   \uncover<1->{Abiertos generadores $A_v$}\\[0.1cm]
   \uncover<2->{\includegraphics[height=0.55cm,angle=45]{Images/Flecha_azul.pdf}\\[0.1cm]
      Vecindades mínimas $U_v=\bigcap_{w\in A_v}A_w$}\\[0.1cm]
   \uncover<3->{\includegraphics[height=0.55cm,angle=45]{Images/Flecha_azul.pdf}\\[0.1cm]
      Base mínima $\mathcal{B}_{\min}=\{U_v \mid v\in V\}$}\\[0.1cm]
   \uncover<4->{\includegraphics[height=0.55cm,angle=45]{Images/Flecha_azul.pdf}\\[0.1cm]
      Uniones arbitrarias}\\[0.1cm]
   \uncover<5->{\includegraphics[height=0.55cm,angle=45]{Images/Flecha_azul.pdf}\\[0.1cm]
      Topología de adyacencia}
\end{frame}

\begin{frame}{Propiedades: topologías discretas}
   \begin{prop}
      La topología $\tau_{\mathcal{A},G}$ es discreta si y solo si $A_x \nsubseteq A_y$ para cualesquiera $x,y\in V$ con $x\neq y$.\vspace{0.3cm}
   \end{prop}
   \vspace{0.3cm}
   \hspace*{-0.22cm}% equilibra las brechas laterales
   \begin{minipage}{\textwidth}
   \begin{overlayarea}{\textwidth}{3.4cm}
   \only<2>{%
   \begin{columns}[c]
      \begin{column}{0.35\textwidth}
         \centering
         \tikzset{every picture/.append style={scale=0.85}}
         \begin{tikzpicture}[
  vnode/.style={circle, fill=FancyBlue!40, draw=black, thick, text=black,
                minimum size=6mm, inner sep=0pt, font=\footnotesize}
]
  % Vértices: v2 en el centro
  \coordinate (v1) at (90:1.25);
  \coordinate (v2) at (0,0);
  \coordinate (v3) at (-45:1.25);
  \coordinate (v4) at (-90:1.25);
  \coordinate (v5) at (225:1.25);

  % Aristas: v1v2, v2v3, v2v4, v2v5
  \foreach \i in {1,3,4,5}
    \draw[thick] (v2) -- (v\i);

  \foreach \i in {1,...,5}
    \node[vnode] at (v\i) {$v_{\i}$};
\end{tikzpicture}

      \end{column}
      \begin{column}{0.63\textwidth}
         $A_{v_1}=\{v_2\}\subseteq\{v_2\}=A_{v_3}$\\[0.3cm]
         Por tanto, su topología de adyacencia no es discreta.
      \end{column}
   \end{columns}}%
   \only<3>{%
   \begin{columns}[c]
      \begin{column}{0.35\textwidth}
         \centering
         \begin{tikzpicture}[
  vnode/.style={circle, fill=FancyBlue!40, draw=black, thick, text=black,
                minimum size=6mm, inner sep=0pt, font=\footnotesize}
]
  % Grafo completo K^3
  \coordinate (v1) at (90:1.0);
  \coordinate (v2) at (210:1.0);
  \coordinate (v3) at (330:1.0);

  \draw[thick] (v1) -- (v2) -- (v3) -- cycle;

  \foreach \i in {1,2,3}
    \node[vnode] at (v\i) {$v_{\i}$};
\end{tikzpicture}

      \end{column}
      \begin{column}{0.63\textwidth}
         $A_{v_1}=\{v_2,v_3\}$\\
         $A_{v_2}=\{v_1,v_3\}$\\
         $A_{v_3}=\{v_1,v_2\}$\\[0.3cm]
         Por tanto, su topología de adyacencia es discreta.
      \end{column}
   \end{columns}}
   \end{overlayarea}
   \end{minipage}
\end{frame}

\begin{frame}{Propiedades: compacidad}
   \begin{prop}
      La topología de adyacencia de un grafo es \alert{compacta} si y solo si el grafo es \alert{finito}.\vspace{0.3cm}
   \end{prop}
\end{frame}

\begin{frame}{Isomorfismo de grafos}
   \begin{itemize}
      \item $G=(V,E)$ y $G'=(V',E')$ dos grafos
      \item $\varphi: V\to V'$ una biyección
      \item $\varphi$ es un \alert{isomorfismo} si y solo si vértices adyacentes de $G$ son adyacentes en $G'$
   \end{itemize}
   \vspace{0.3cm}
   \begin{center}
      \begin{tikzpicture}[
  scale=0.4,
  vnode/.style={circle, fill=FancyBlue!40, draw=black, thick, text=black,
                minimum size=6mm, inner sep=0pt, font=\footnotesize},
  vino/.style={vnode, fill=Vino!45}
]
  \definecolor{Vino}{HTML}{503642}% vino tinto del póster
  % Grafos isomorfos G y G' (Figura del documento: G y G' son isomorfos)
  % ---------- G ----------
  \begin{scope}[xshift=-5.5cm]
    \node[vino] (v4) at (-3, 3.5) {$v_4$};
    \node[vino] (v5) at (3, 3.5) {$v_5$};
    \node[vino] (v6) at (3, -2.5) {$v_6$};
    \node[vino] (v7) at (-3, -2.5) {$v_7$};
    \node[vnode] (v1) at (-1.5, 1.5) {$v_1$};
    \node[vnode] (v2) at (1.5, 1.5) {$v_2$};
    \node[vnode] (v3) at (0, -0.5) {$v_3$};
    \foreach \from/\to in {v1/v2, v2/v3, v3/v1, v3/v7, v3/v6}
      \draw[thick] (\from) -- (\to);
    \foreach \from/\to in {v4/v5, v5/v6, v6/v7, v7/v4}
      \draw[thick, Vino] (\from) -- (\to);
    \node at (0,-4.3) {$G$};
  \end{scope}
  % ---------- G' ----------
  \begin{scope}[xshift=5.5cm]
    \node[vnode] (w3) at (0, 3.5) {$w_3$};
    \node[vino] (w7) at (-1.2, 1.5) {$w_7$};
    \node[vino] (w6) at (1.2, 1.5) {$w_6$};
    \node[vino] (w5) at (1.2, -0.5) {$w_5$};
    \node[vino] (w4) at (-1.2, -0.5) {$w_4$};
    \node[vnode] (w1) at (-3, -2.5) {$w_1$};
    \node[vnode] (w2) at (3, -2.5) {$w_2$};
    \foreach \from/\to in {w3/w7, w3/w6, w1/w2}
      \draw[thick] (\from) -- (\to);
    \foreach \from/\to in {w7/w6, w6/w5, w5/w4, w4/w7}
      \draw[thick, Vino] (\from) -- (\to);
    \draw[thick] (w3) to[out=180, in=90] (w1);
    \draw[thick] (w3) to[out=0, in=90] (w2);
    \node at (0,-4.3) {$G'$};
  \end{scope}
\end{tikzpicture}

   \end{center}
\end{frame}

\begin{frame}{Propiedades: isomorfismo de grafos}
   \begin{prop}
      Si dos grafos son \alert{isomorfos}, sus respectivos espacios topológicos son \alert{homeomorfos}.\vspace{0.3cm}
   \end{prop}
\end{frame}

\begin{frame}{Propiedades: isomorfismo de grafos}
   \begin{Obs}
      El recíproco no es cierto.\vspace{0.3cm}
   \end{Obs}
   \vspace{0.4cm}
   \noindent\hspace*{10.7pt}\alert{Ejemplo}
   \begin{center}
      \begin{tikzpicture}[
  vnode/.style={circle, fill=FancyBlue!40, draw=black, thick, text=black,
                minimum size=5mm, inner sep=0pt, font=\scriptsize}
]
  % K^5 y C^5: espacios de adyacencia homeomorfos (ambos discretos), grafos no isomorfos
  % ---------- K^5 ----------
  \begin{scope}[xshift=-2.2cm]
    \foreach \i in {1,...,5}
      \coordinate (k\i) at ({90+72*(\i-1)}:0.8);
    \foreach \i/\j in {1/2,1/3,1/4,1/5,2/3,2/4,2/5,3/4,3/5,4/5}
      \draw[thick] (k\i) -- (k\j);
    \foreach \i in {1,...,5}
      \node[vnode] at (k\i) {$v_{\i}$};
    \node at (0,-1.2) {$K^5$};
  \end{scope}
  % ---------- C^5 ----------
  \begin{scope}[xshift=2.2cm]
    \foreach \i in {1,...,5}
      \coordinate (c\i) at ({90+72*(\i-1)}:0.8);
    \draw[thick] (c1) -- (c2) -- (c3) -- (c4) -- (c5) -- cycle;
    \foreach \i in {1,...,5}
      \node[vnode] at (c\i) {$w_{\i}$};
    \node at (0,-1.2) {$C^5$};
  \end{scope}
\end{tikzpicture}

   \end{center}
\end{frame}
